% Geometría conceptual: la figura ES la tesis (aquí un cono de secciones). Solo cabe cuando
% el texto usa o justifica la figura. Ejes globales (contracción/dilatación) al margen.
\documentclass[tikz,border=12pt]{standalone}
\usepackage{diagrama}
\begin{document}
\begin{tikzpicture}[every node/.append style={font=\px{28}}]
  \node[titulo=44] at (3,7.2) {Título de la figura};
  \draw[puntos] (1,6.2) -- (3,0.8);                  % cono: lados punteados
  \draw[puntos] (5,6.2) -- (3,0.8);
  \draw[rayas] (3,6.2) ellipse [x radius=2, y radius=0.55];   % secciones: contorno virtual
  \draw[rayas] (3,4.2) ellipse [x radius=1.3, y radius=0.38];
  \draw[rayas] (3,2.4) ellipse [x radius=0.6, y radius=0.18];
  \node (n1) at (3,6.2) {S1}; \node (n2) at (3,4.2) {S2}; \node (n3) at (3,2.4) {S3};
  \node[font=\bfseries\px{34}] (n4) at (3,0.3) {Presente};
  \draw[eje] (-0.4,5.8) -- node[rotulo=20, rotate=90] {\textit{Contracción}} (-0.4,1.2);  % x1
  \draw[eje] (6.4,1.2) -- node[rotulo=20, rotate=90] {\textit{Dilatación}} (6.4,5.8);     % x2
\end{tikzpicture}
\end{document}
